\chapter*{Notación, tipos y desambiguación de símbolos}
\addtocontents{toc}{}
\addcontentsline{toc}{chapter}{Notación, tipos y desambiguación de símbolos}

\begin{longtable}{@{}p{.22\textwidth}p{.70\textwidth}@{}}
\toprule
Símbolo & Objeto \\ \midrule
\endhead
\(\mathcal A_9\)&Toro aritmético nonádico de caminos con dos evaluaciones
internas.\\
\(\mathrm{TRIT}\)&Coordenada ternaria orientada del transporte; la salida
visible no agota el levantamiento enriquecido con acarreo.\\
\(\mathrm{TPK}\)&Núcleo topológico de fase con estado enriquecido,
transiciones, rutas, prolongaciones y memoria.\\
\(M_{\rm ph}\)&Selector ternario de fase operativa, distinto del transporte
cardinal y del muestreo estroboscópico.\\
\(\tau_3\)&Bloque primitivo de fase \((+1,-1,0)\).\\
\(\mathcal N,\mathcal E,\mathcal S,\mathcal O\)&Generadores del transporte
cardinal.\\
\(\operatorname{Obs}_3\)&Muestreo estroboscópico de cada tercer paso.\\
\(\mathcal R_{12}\)&Sistema temporal dodecafásico orientado interno a TPK.\\
\(\Theta_{108}=\mathbb Z/108\mathbb Z\)&Calendario observable
\(12\times9\); no es el estado TPK completo.\\
\(C_{12}\)&Rotación de los doce sectores temporales.\\
\(\Gamma_{108}\)&Órbita orientada de transporte de \(108\) pasos.\\
\(\mathcal G_9\)&Poda unitaria dentro de la relación de extensión y
frontera; no constituye por sí sola la monodromía.\\
\(\mathfrak M_9\)&Monodromía nonádica con memoria: retorna la fase y cambia
el estado.\\
\(\pi_{\rm term}\)&Proyección del estado terminal enriquecido a los tres
bloques visibles seleccionados.\\
\(R_{\rm sgn}\)&Lectura de la coordenada entera firmada del mismo estado
terminal enriquecido.\\
\(U_{12}^{\rm cat}\)&Concatenación visible \(U_6\Vert U_6\).\\
\(U_{12}^{\rm sgn}\)&Coordenada firmada del estado enriquecido, reversible
con \(K\) mediante \(\mathcal H_{12}\); no se deduce de
\(U_{12}^{\rm cat}\), de la proyección calendárica ni de los tres bloques terminales
aislados.\\
\(\mathcal H_{12}\)&Transformación dodecafásica reversible entre
\(U_{12}^{\rm sgn}\) y \(K\).\\
\(\mathcal E_{\rm dod}\)&Extractor integral
\((D_3,D_4,Q):\mathbb Z^{12}\to\mathbb Z^{25}\) que reconstruye \(K\).\\
\(\overline{\mathcal C}_9^{\,3}\)&Completación nonádica cúbica.\\
\(A,C^*\)&Coordenadas angular y contraangular intrínsecas.\\
\(q_\pm\)&Canales orientados del lector bicapa.\\
\(D_{\rm raw}^{90,120}\)&Diferencia diagnóstica cruda de series en una
carta fijada; no es un extractor canónico.\\
\(\iota_{6,8}\)&Inclusión
\(\mathcal F_6\hookrightarrow\mathcal F_8\) de incidencias hexada--octada.\\
\(K_{6\to8}\)&Funcional angular de las dos subcolas de Lambert, posterior a
\(q_\pm\), las series, el polo, la positividad y la minimalidad.\\
\(\langle90,120\rangle\)&Semigrupo espectral que gobierna la jerarquía de
torres.\\
\(\varepsilon_{\rm res}^{\circ}\)&Residuo de la identidad angular exacta;
no es una extracción finita de los objetos 90/120 anteriores.\\
\(\Delta_4\)&Torsión mínima o invariante torsional global anterior a la
especie; no es la torsión de Cartan.\\
\(\mathsf T_{\alpha,\infty}\)&Prolongación remota del acarreo de la vía
entera de \(\alpha\), tipada separadamente del generador orbital
\(T_1^{\rm orb}=L_1\), del contador neutral y del calendario.\\
\(N_1^{\rm neu}\)&Número de posiciones neutrales en la primera ventana del
ejemplo de transporte; vale $3$ y no es un operador.\\
\(T_1^{\rm orb}:=L_1\)&Segundo generador en el calibre orbital local; en su
sección de referencia se abrevia \(T_1\).\\
\(\varepsilon_{\rm clk}\)&Cuanto energético del reloj,
\(\hbar_{\rm int}2\pi/(108t_0)\); no es una base energética canónica sin
declarar el marco temporal. El símbolo \(\varepsilon_0\) queda reservado a
la permitividad del vacío en la interfaz constitutiva SI.\\
\(\mathrm{CLK}1\text{--}\mathrm{CLK}6\)&Etiquetas de las seis ecuaciones del
reloj finito, desde el Hamiltoniano hasta Schrödinger.\\
\(\mathcal K_{27}\)&Retorno compuesto de veintisiete actualizaciones.\\
\(X_\infty^{\rm cal}\)&Dominio enriquecido común de las evaluaciones
globales.\\
\bottomrule
\end{longtable}

\section*{Cardinales repetidos con fuentes diferentes}

La igualdad de un cardinal no transporta automáticamente dominio,
operación, memoria ni fuerza probatoria. La tabla siguiente fija las
colisiones que más afectan a la lectura.

\begin{longtable}{@{}p{.12\textwidth}p{.27\textwidth}p{.51\textwidth}@{}}
\toprule
Índice & Objeto tipado & Función y separación \\ \midrule
\endhead
\(108\)&\(\Theta_{108}\)&Calendario \(12\times9\); cierra el reloj
observable.\\
\(108\)&\(\Pi_{108}^{\rm cal}\)&Proyección calendárica de un cursor con \(324\)
semillas; olvida segundo cursor, orientación, acarreo, frontera, registro
cronológico enriquecido y
supervivencia.\\
\(108\)&\(w_{108}\)&Palabra de \(108=18\times6\) trits; no es el calendario.\\
\(104,976\)&\(X_{\rm TPK}^{\rm seed}\)&Motor de dos cursores,
\(104,976=324^2\); conserva un dominio distinto del calendario.\\
\(1080\)&\(w_{1080}\)&Longitud de una palabra.\\
\(1080\)&\(\Gamma_0(1080)\)&Nivel de un subgrupo modular; requiere su
realización propia.\\
\(729\)&\(\mathbb F_3^6\)&Ambiente de palabras ternarias.\\
\(729\)&\(9^3\)&Célula cúbica/interior EMRH.\\
\(729\)&\(\operatorname{Ext}(X_K)\)&Número de subcilindros examinados en cada
etapa de prolongación.\\
\(729\)&bulk EMRH&Estrato interior de la completación base mil.\\
\(27\)&\(\mathcal K_{27}\)&Operador de retorno compuesto; no nombre de una
teoría autónoma.\\
\(27\)&estrato cúbico&Contribución de dos coordenadas de frontera en
\(999=729+243+27\).\\
\(27\)&módulo o clasificador&Cada aparición exige su mapa antes de
identificarse con otra.\\
\bottomrule
\end{longtable}

\section*{Seis fuentes distintos de 90/120}

\begin{longtable}{@{}
>{\raggedright\arraybackslash}p{.21\textwidth}
>{\raggedright\arraybackslash}p{.42\textwidth}
>{\raggedright\arraybackslash}p{.27\textwidth}@{}}
\toprule
Objeto & Dominio y operación & Salida y prohibición de tipo \\ \midrule
\endhead
\(\mathcal E_{\rm dod}\)&\(K\in\mathbb Z^{12}\); diferencias por pasos
\(3,4\) y suma \(Q\)&Coordenadas transversales dodecafásicas; no incidencia
ni serie angular.\\
\(D_{\rm raw}^{90,120}\)&Una carta \(q\); diferencia cruda
\(S_{90}(q)-S_{120}(q)\)&Diagnóstico de escala, distinto del funcional de
Barbero--Immirzi y de los operadores de incidencia.\\
\(\iota_{6,8}\)&\(H\subset O\); inclusión de
\(H\times\binom{H}{2}\) en \(O\times\binom{H}{2}\)&Cardinales \(90,120\) y
defecto \(-1/12\); no extractor ni serie.\\
\(K_{6\to8}\)&Canales \(q_\pm\) y subcolas de Lambert; combinación
\(12\Delta S_{90}-\Delta S_{120}\)&Funcional angular relativo a sus
premisas; no se transporta como bloque de masa.\\
\(\mathfrak S=\langle90,120\rangle\)&Alturas de momentos
\(s_n=q_-^n-q_+^n\)&Semigrupo global de torres; no termina en ocho alturas.\\
\bottomrule
\end{longtable}

\section*{Cuatro realizaciones de la reversibilidad bidireccional}

\begin{longtable}{@{}
  >{\raggedright\arraybackslash}p{.26\textwidth}
  >{\raggedright\arraybackslash}p{.26\textwidth}
  >{\raggedright\arraybackslash}p{.38\textwidth}@{}}
\toprule
Operación & Dominio & Acción \\ \midrule
\endhead
\(P_\pm=(I\pm R)/2\)&Álgebra escindida&Proyecta los dos sectores
espectrales.\\
Inversión de rumbo&Ruta y calendario compuesto&Invierte la orientación del
transporte sin borrar su palabra.\\
\(\operatorname{Ext}^{\pm}\)&Árbol de prolongaciones&Compara valencias
futuras y pasadas; su balance puede leerse mediante
\(\log(\#\operatorname{Ext}^+/\#\operatorname{Ext}^-)\).\\
\(C_M(\tau)=-\operatorname{rev}(\tau)\)&Palabra dodecafásica&Conjuga orden
y signo; no es la involución celular \(\rho_c\).\\
\bottomrule
\end{longtable}

La función común de estas cuatro realizaciones es conservar orientación
reversible. No autoriza a identificar retorno de fase, retorno de ruta,
retorno de palabra y retorno de estado.
